% ---------------------------------
\section{Results}

\subsection{Research Findings (Perplexity
Output)}

Perplexity returned structured research
across more than 10 verified resources. 
The main privacy risks in modern AI systems were
data minimization failures, purpose limitation violations,
consent bypassing, and unauthorized secondary use of training data.
These risks are intensified by weak consent mechanisms, 
limited transparency, insecure data-sharing pipelines, and 
third-party vendors.
Privacy by Design, differential privacy, federated learning,
access controls, and continuous audits were identified as 
important safeguards.
These key findings of the Perplexity enable the Claude AI to create 
a comprehensive document to be used by the other AI tools like 
QuillBot and SlidesAI. 

---

\subsection{Synthesized Content (Claude API
Output)}

Claude AI created a detailed document to be used 
in other AI tools. However, the statement below 
probably hallucinated the need to create a “Microsoft Word” 
document, which was not intended:

\begin{framed}
\small
\texttt{``write a 1500-2000 word
article on Privacy and Data Protection
in AI Systems''}
\end{framed}


For demonstration and research purposes, 
it was better to keep this error.

The problem with this kind of document generation 
is the number of tokens spent in other AI tools, 
such as ChatGPT. These tools need to parse the 
whole document to extract the relevant content, 
which consumes more tokens than necessary. 
A better format would be Markdown. 
Therefore, next time, the prompt should be more 
specific about the expected output format.
---

\subsection{Writing Style Comparison
(QuillBot Output)}

To see how writing style affects clarity,
one paragraph from the Claude output was run
through QuillBot. Only the Standard mode was
available, since the free tier locks Formal
and Creative behind a subscription, so the
comparison is limited to the original
against a single rewrite.

---

\subsection{Presentation Output
(SlidesAI.io)}

For the deck, ChatGPT was used first as a
meta-prompting step: GPT-5.5 generated the
SlidesAI prompt and suggested templates, of
which the Soft Sage one was chosen, and that
prompt was then handed to SlidesAI.io.
SlidesAI produced a 10-slide presentation
covering the title, context, the main
problem, key concepts, the approach, the
main findings, supporting examples from the
source, implications, practical takeaways,
and a conclusion. It ran in about ten
minutes, selected typography and images on
its own, and needed no manual editing, which
makes it suitable for classroom use or wider
dissemination.

---
